\documentclass[12pt]{article}

% ── Packages ──────────────────────────────────────────────────────────────────
\usepackage[margin=1in]{geometry}
\usepackage{amsmath, amssymb, amsthm}
\usepackage{mathtools}
\usepackage[T1]{fontenc}
\usepackage[utf8]{inputenc}
\usepackage{hyperref}
\usepackage{xcolor}
\usepackage{mdframed}
\usepackage{enumitem}
\usepackage{titlesec}
\usepackage{booktabs}
\usepackage{setspace}
\usepackage{parskip}

% ── Colors ────────────────────────────────────────────────────────────────────
\definecolor{lawblue}{RGB}{15, 52, 96}
\definecolor{conjectpurple}{RGB}{123, 45, 139}
\definecolor{titlenavy}{RGB}{26, 26, 46}
\definecolor{darkgray}{RGB}{60, 60, 60}

% ── Hyperref ──────────────────────────────────────────────────────────────────
\hypersetup{
  colorlinks=true,
  linkcolor=lawblue,
  citecolor=lawblue,
  urlcolor=lawblue
}

% ── Section formatting ────────────────────────────────────────────────────────
\titleformat{\section}{\large\bfseries\color{titlenavy}}{}{0em}{}[\titlerule]
\titleformat{\subsection}{\normalsize\bfseries\color{lawblue}}{}{0em}{}
\titleformat{\subsubsection}{\normalsize\bfseries\color{darkgray}}{}{0em}{}
\titlespacing{\section}{0pt}{18pt}{8pt}
\titlespacing{\subsection}{0pt}{14pt}{6pt}
\titlespacing{\subsubsection}{0pt}{10pt}{4pt}

% ── Theorem environments ──────────────────────────────────────────────────────
\newtheoremstyle{lawstyle}
  {10pt}{10pt}{\itshape}{0pt}{\bfseries\color{lawblue}}{.}{.5em}{}

\newtheoremstyle{conjstyle}
  {10pt}{10pt}{\itshape}{0pt}{\bfseries\color{conjectpurple}}{.}{.5em}{}

\newtheoremstyle{predstyle}
  {8pt}{8pt}{\normalfont}{0pt}{\bfseries}{.}{.5em}{}

\theoremstyle{lawstyle}
\newtheorem{law}{Law}

\theoremstyle{conjstyle}
\newtheorem{conjecture}{Conjecture}[section]

\theoremstyle{predstyle}
\newtheorem{prediction}{Prediction}

% ── Framed environments ───────────────────────────────────────────────────────
\newmdenv[
  linecolor=lawblue,
  linewidth=2pt,
  topline=false, rightline=false, bottomline=false,
  leftmargin=12pt, innerleftmargin=10pt,
  rightmargin=0pt, innerrightmargin=0pt,
  innertopmargin=4pt, innerbottommargin=4pt,
  skipabove=8pt, skipbelow=8pt
]{lawbox}

\newmdenv[
  linecolor=conjectpurple,
  linewidth=2pt,
  topline=false, rightline=false, bottomline=false,
  leftmargin=12pt, innerleftmargin=10pt,
  rightmargin=0pt, innerrightmargin=0pt,
  innertopmargin=4pt, innerbottommargin=4pt,
  skipabove=8pt, skipbelow=8pt
]{conjbox}

% ── Title ─────────────────────────────────────────────────────────────────────
\title{
  {\Large\bfseries\color{titlenavy} Physical Mathematics: A Foundational Framework}
}
\author{}
\date{June 2025}

% ─────────────────────────────────────────────────────────────────────────────
\begin{document}
\maketitle
\thispagestyle{empty}

% ── Abstract ──────────────────────────────────────────────────────────────────
\begin{abstract}
Anson's Framework for Physical Mathematics proposes that mathematical truth is not timeless
and substrate-independent, but physically realized---constrained by thermodynamics, reference
frames, quantum mechanics, fractal geometry, and spatial structure. The framework consists of
\textbf{two laws} and \textbf{three conjectures}, distinguished by their current epistemic
status but held to identical standards of formal rigor. The ideas, conceptual architecture,
and epistemological positions presented here are original. The mathematical formalization was
developed through human-AI coupling, as disclosed in the Acknowledgments.

\textbf{Anson's First Law} (Thermodynamic Constraint) derives a computable precision ceiling
$n_{\max}$ for any physical system from Landauer's Principle and the Bekenstein bound.
\textbf{Anson's Second Law} (Frame Dependence) proves three mechanically verified properties
of realized computation: (i) both error terms are strictly positive in any finite frame;
(ii) truncation and rounding errors are strictly and independently monotone in opposite
directions; (iii) the optimal step size is uniquely and positively determined by the physical
frame. The law does not assert a strict total ordering of error between frames---that claim
is unprovable, as the terms can cancel---but proves the correct, stronger claim that each
term is independently frame-determined.
\textbf{Conjecture~A} (Quantum Complexity Reduction) proposes that complexity class membership
is substrate-dependent---a structural claim with a specific falsification condition.
\textbf{Conjecture~B} (Fractal Dimensional Correspondence) proposes that arithmetic
irregularity at finite resolution is a projection artifact of underlying fractal symmetry,
with a computable convergence test. \textbf{Conjecture~C} (Geometric Foundation) proposes
that numerical systems are coordinate representations of prior geometric structures, with a
spectral gap falsification criterion. The two laws are formally verified in Lean~4~\cite{lean4} (Mathlib~\cite{mathlib2020});
 the three research claims are recorded as unasserted propositions with precise type signatures. All five components
generate experimentally testable predictions.
\end{abstract}

\bigskip
\hrule
\bigskip

% ── 1. Introduction ───────────────────────────────────────────────────────────
\section{Introduction}

Standard mathematics treats its objects as substrate-independent: the Riemann zeros, the
primes, NP certificates exist timelessly, and computation is merely a window onto
pre-existing truth. This framework rejects that position. The act of mathematical
verification is irreducibly physical---it costs energy, occupies space, operates in a
reference frame, and is bounded by the same laws that govern any physical process.
Mathematical truth, on this account, is what can be \emph{realized}, not merely what can
be stated.

This position has antecedents. Landauer~\cite{landauer1961} established thermodynamic cost
for logical irreversibility; Bekenstein~\cite{bekenstein1981} derived a universal entropy
bound; Deutsch~\cite{deutsch1985} argued that physics determines computability;
Wheeler~\cite{wheeler1990} proposed information as ontologically fundamental. The present
framework extends this lineage into a unified formal structure with explicit falsification
conditions.

\subsection*{On the distinction between Laws and Conjectures}

A \textbf{law} in this framework is a statement derivable from established physical
principles (Landauer, Bekenstein, numerical analysis) whose falsification conditions follow
from those same principles. A \textbf{conjecture} is a statement that is formally precise,
internally consistent, and falsifiable, but whose derivation from first principles is
incomplete---it requires either experimental confirmation or a proof step not yet supplied.
The conjectures are held to exactly the same standard of formal rigor as the laws. The
distinction is epistemic, not ontological: conjectures are laws in waiting, not speculations.

\bigskip
\hrule
\bigskip

% ── 2. The Two Laws ───────────────────────────────────────────────────────────
\section{The Two Laws}

% ── Law 1 ────────────────────────────────────────────────────────────────────
\subsection{Anson's First Law: Thermodynamic Constraint on Computational Precision}

\subsubsection*{Statement}

\begin{law}[Thermodynamic Constraint]
For any physical system with energy bound $E_{\max}$ and spatial radius $R$, the maximum
number of bits that can be reliably verified is:
\begin{lawbox}
\[
  n_{\max} = \left\lfloor \frac{2\pi k_B R\, E_{\max}}{\hbar\, c\, \ln 2} \right\rfloor
\]
\end{lawbox}
Any mathematical computation requiring verification to precision $b(n) > n_{\max}$ bits is
physically unrealizable within that system. This is an absolute constraint, not an
engineering limitation.
\end{law}

\subsubsection*{Derivation}

By Landauer's Principle~\cite{landauer1961}, each irreversible bit operation dissipates at
least $E_{\mathrm{bit}} = k_B T \ln 2$. Verifying a mathematical statement to $b$ bits
requires at least $b$ irreversible operations:
\[
  E_{\mathrm{verify}}(b) \;\geq\; b \cdot k_B T \ln 2.
\]
The Bekenstein bound~\cite{bekenstein1981} constrains the total entropy of any bounded
physical system:
\[
  S_{\max} = \frac{2\pi k_B R E}{\hbar c}.
\]
Each verified bit contributes $k_B \ln 2$ to entropy, so:
\[
  b \;\leq\; \frac{S_{\max}}{k_B \ln 2} = \frac{2\pi R E}{\hbar c \ln 2}.
\]
Setting $E = E_{\max}$ yields $n_{\max}$. The two bounds are both necessary: Landauer sets
the per-bit cost; Bekenstein sets the total budget. This is not a statement about computation
speed---it is a statement about how much mathematical truth a bounded physical system can
simultaneously hold.

\subsubsection*{Falsifiable Prediction}

\begin{prediction}[P1]
Energy consumed verifying $n$ Riemann zeros (via \texttt{mpmath}, GPU power-monitored) must
scale as $E(n) \geq \alpha \cdot n \log n$ with $\alpha \geq k_B T \ln 2 \cdot c_0$, where
$c_0$ is the known bit-cost of the algorithm. Any measured $\alpha$ below this floor falsifies
the First Law.
\end{prediction}

\subsubsection*{Scope and Limitations}

The bound is tight only at the Landauer limit. Room-temperature hardware operates orders of
magnitude above this; $n_{\max}$ is an idealized ceiling. The law is silent on \emph{logical}
unprovability (G\"odel)---it addresses \emph{physical} unrealizability. These are distinct
constraints that may independently apply to the same statement.

\bigskip

% ── Law 2 ────────────────────────────────────────────────────────────────────
\subsection{Anson's Second Law: Computational Frame Dependence}

\subsubsection*{Statement}

\begin{law}[Frame Dependence]
Let $T_{\mathrm{ideal}}$ be the exact solution to a computation $C$. In a physical frame
characterized by precision $p$ (bits), time budget $\tau$, and numerical method of local
truncation order $q$ with step size $h = \tau/N$, three properties hold:

\begin{lawbox}
\textbf{(i) Positivity.} Both error terms are strictly positive in any finite physical frame:
\[
  C_{\mathrm{method}} \cdot h^q > 0 \qquad \text{and} \qquad C_{\mathrm{round}} \cdot p^{-1} > 0.
\]
\textbf{(ii) Tradeoff.} The terms are independently monotone in opposite directions: increasing
$h$ strictly increases truncation error; increasing $p$ strictly decreases rounding error:
\[
  h_1 < h_2 \;\Rightarrow\; C_{\mathrm{method}} \cdot h_1^q < C_{\mathrm{method}} \cdot h_2^q,
\]
\[
  p_1 < p_2 \;\Rightarrow\; C_{\mathrm{round}} \cdot p_2^{-1} < C_{\mathrm{round}} \cdot p_1^{-1}.
\]
\textbf{(iii) Optimal frame.} The step size minimizing total error is uniquely determined by the
physical frame parameters:
\[
  h^* = \left(\frac{C_{\mathrm{round}}}{q\, C_{\mathrm{method}}\, p}\right)^{1/(q+1)} > 0.
\]
\end{lawbox}

The frame $(p, \tau, q)$ is thermodynamically bounded: $p$ by hardware word length, $\tau$ by
energy budget via the First Law, $q$ by memory. Because the two error terms pull in opposite
directions and neither can reach zero, frame dependence is irreducible---not merely practical.

\textbf{Note on total error ordering.} The total error $C_{\mathrm{method}} h^q + C_{\mathrm{round}} p^{-1}$
is \emph{not} in general strictly ordered between frames: the truncation and rounding terms can
cancel, making two distinct frames yield equal total error. What is provable---and proved---is
that the individual terms are strictly and independently frame-dependent. This is the correct
and precise content of the law.
\end{law}

\subsubsection*{Derivation}

The error decomposition follows from standard numerical analysis~\cite{higham2002}. For a
$q$-th order method, global discretization error is $O(h^q)$, with constant $C_{\mathrm{method}} > 0$.
Floating-point arithmetic at $p$ bits introduces rounding error bounded by $C_{\mathrm{round}} \cdot p^{-1}$,
where $p^{-1} > 0$ for any finite $p$.

\textit{Proof of positivity (i).} $C_{\mathrm{method}} > 0$ and $h^q > 0$ for $h > 0$, so
$C_{\mathrm{method}} \cdot h^q > 0$. Similarly $C_{\mathrm{round}} > 0$ and $p^{-1} > 0$ for
$p > 0$, so $C_{\mathrm{round}} \cdot p^{-1} > 0$. Both hold for all finite physical frames. $\square$

\textit{Proof of tradeoff (ii).} For fixed positive $C_{\mathrm{method}}$, the map $h \mapsto h^q$
is strictly increasing for $q \geq 1$, so $h_1 < h_2$ implies $h_1^q < h_2^q$, giving strict
truncation ordering by multiplicativity. For the rounding term: $p_1 < p_2$ with both positive
implies $p_2^{-1} < p_1^{-1}$ (inversion reverses order), so $C_{\mathrm{round}} \cdot p_2^{-1} <
C_{\mathrm{round}} \cdot p_1^{-1}$ by multiplicativity. $\square$

\textit{Proof of optimal step (iii).} Minimize $f(h) = C_{\mathrm{method}} h^q + C_{\mathrm{round}} p^{-1} h^{-1}$
over $h > 0$. Setting $f'(h) = 0$: $q\, C_{\mathrm{method}} h^{q-1} = C_{\mathrm{round}} p^{-1} h^{-2}$,
giving $h^* = (C_{\mathrm{round}} / (q\, C_{\mathrm{method}}\, p))^{1/(q+1)} > 0$. $\square$

The physical grounding: the First Law bounds $\tau$, which bounds $N = \tau/h^*$. The two
laws are coupled---the First Law constrains the parameter space in which the Second Law operates.

\subsubsection*{Lean 4 Verification}

The three properties above are mechanically verified in Lean~4~\cite{lean4} (Mathlib~\cite{mathlib2020}) as:
\begin{itemize}
  \item \texttt{law2\_error\_bound\_pos}: positivity of both terms and their sum.
  \item \texttt{law2\_truncation\_monotone}, \texttt{law2\_rounding\_monotone}: strict
        independent monotonicity of each term.
  \item \texttt{law2\_tradeoff}: the conjunction of both monotonicity results.
  \item \texttt{optimalStep\_pos}: the optimal step is well-defined and strictly positive.
\end{itemize}
The Lean proof process identified that a stronger claim---strict ordering of \emph{total}
error between frames---is unprovable, since the terms can cancel. The law as stated above
reflects exactly what is mechanically verified.

\subsubsection*{Falsifiable Predictions}

\begin{prediction}[P2a]
Lorenz system solved with RK4 at $p \in \{32, 64, 128\}$ bits, fixed $\tau$, pre-chaos
regime ($t < t_{\mathrm{Lyapunov}}$): truncation error must increase strictly with $h$
and rounding error must decrease strictly with $p$, independently and simultaneously.
If either monotonicity fails, the tradeoff claim is falsified.
\end{prediction}

\begin{prediction}[P2b]
The empirically optimal step $h^*$ at each $p$ must match the analytic minimizer
$(C_{\mathrm{round}} / (q\, C_{\mathrm{method}}\, p))^{1/(q+1)}$ to within experimental
precision. The optimal step must be strictly positive and vary monotonically with $p$.
\end{prediction}

\subsubsection*{Scope and Limitations}

The error decomposition is standard numerical analysis. The law's claim is that \emph{physical}
constraints on $(p, \tau)$ make both error terms irreducibly frame-dependent. The law does not
claim that different frames always produce different total errors---it claims that the individual
truncation and rounding contributions are each strictly and independently determined by the
physical frame. This is the weaker, correct, and fully verified claim.

\bigskip
\hrule
\bigskip

% ── 3. The Three Conjectures ──────────────────────────────────────────────────
\section{The Three Conjectures}

The following three conjectures are formally specified to the same standard as the laws
above. Each carries a precise statement, a derivation sketch identifying the missing step,
and two independent falsification paths. They are called conjectures because the derivation
from first principles is incomplete---not because the claims are imprecise or untestable.

% ── Conjecture A ──────────────────────────────────────────────────────────────
\subsection{Conjecture A: Quantum Complexity Reduction}

\subsubsection*{Statement}

\begin{conjecture}[Quantum Complexity Reduction]
Complexity class membership is substrate-dependent. For any problem $\Pi \in \mathbf{NP}$
with classical worst-case complexity $C_{\mathrm{cl}}(n)$, the complexity $C_q(n)$ realized
by a quantum substrate satisfies $C_q(n) = o(C_{\mathrm{cl}}(n))$ for at least a proper
subclass of $\mathbf{NP}$. Classical and quantum systems are distinct physical frames in the
sense of the Second Law, and the complexity experienced by a computation is frame-relative.

\begin{conjbox}
Canonical instances:
\begin{align*}
  \text{Factoring:} &\quad C_{\mathrm{cl}} \sim \exp\!\left(O(n^{1/3})\right)
    \;\longrightarrow\; C_q \sim O(n^2 \log n) \quad \text{\cite{shor1997}}\\
  \text{Unstructured search:} &\quad C_{\mathrm{cl}} \sim O(N)
    \;\longrightarrow\; C_q \sim O(\sqrt{N}) \quad \text{\cite{grover1996}}
\end{align*}
\end{conjbox}
\end{conjecture}

\subsubsection*{Derivation Sketch and Missing Step}

Shor's algorithm operates in a $2^n$-dimensional Hilbert space, using quantum Fourier
transform to find periods of modular exponentiation in $O(n^2 \log n)$ gate operations. The
classical sub-exponential barrier arises because classical registers cannot coherently
superpose $2^n$ states---a physical constraint (decoherence), not a logical one.

\textbf{Missing step:} A proof that this reduction is \emph{substrate-general}---i.e., that
for any quantum-realizable physical system the complexity gap is preserved---requires a
formal characterization of which physical substrates instantiate $\mathbf{BQP}$ and which do
not. Current results establish the gap for specific algorithms; the conjecture asserts it is
a structural feature of substrate physics.

\subsubsection*{Falsifiable Prediction}

\begin{prediction}[PA]
Grover 3-SAT over $N = 2^{20}$ instances on a 20-qubit processor: runtime must scale as
$O(\sqrt{N})$ vs.\ $O(N)$ classically. The crossover as a function of decoherence rate
$\gamma$ provides a joint test; if runtime scales as $O(\sqrt{N} \cdot e^{\gamma t})$ for
measurable $\gamma$, the substrate-frame coupling is empirically confirmed.
\end{prediction}

\subsubsection*{What Would Falsify It}

If a quantum substrate correctly implements Grover's algorithm at the gate level yet runtime
scales as $O(N)$, and this cannot be explained by decoherence, Conjecture~A is falsified for
that substrate---requiring revision of the substrate-dependence claim.

\bigskip

% ── Conjecture B ──────────────────────────────────────────────────────────────
\subsection{Conjecture B: Fractal Dimensional Correspondence}

\subsubsection*{Statement}

\begin{conjecture}[Fractal Dimensional Correspondence]
Let $S_n = \{s_1, \ldots, s_n\}$ be a finite sample of a number-theoretic sequence (e.g.,
imaginary parts of Riemann zeros). Let $D_H(n)$ be the box-counting Hausdorff dimension of
$S_n$. Then: (i) the limit $D_H = \lim_{n\to\infty} D_H(n)$ exists; (ii) $D_H \in (1,2)$;
and (iii) the apparent irregularity of $S_n$ at finite $n$ is a projection artifact of
finite coarse-graining resolution, not an intrinsic property of the sequence.

Define coarse-graining $P_\varepsilon$ by covering $S_n$ with boxes of size $\varepsilon$
and retaining box-centers. Then:
\begin{conjbox}
\[
  P_\varepsilon(S_\infty) \;\cong\; S_n \quad \text{as } \varepsilon \to 1/n,
\]
with box count satisfying $N(\varepsilon) \sim \varepsilon^{-D_H}$ in the scaling regime.
\end{conjbox}
\end{conjecture}

\subsubsection*{Derivation Sketch and Missing Step}

Montgomery~\cite{montgomery1973} and Odlyzko~\cite{odlyzko1987} showed that spacings between
Riemann zeros follow GUE (Gaussian Unitary Ensemble) statistics. GUE spacing distributions
are scale-invariant in the bulk---a necessary condition for fractal structure.
Berry and Keating~\cite{berry1999} connected Riemann zeros to the spectrum of a hypothetical
quantum Hamiltonian, reinforcing the scale-invariance link.

\textbf{Missing step:} Scale-invariance of GUE statistics is necessary but not sufficient
for a well-defined Hausdorff dimension. The conjecture requires showing that the
box-counting limit converges, and that $D_H$ is universal across the GUE family---not an
artifact of the zero-counting measure.

\subsubsection*{Falsifiable Prediction}

\begin{prediction}[PB]
Box-counting of the first $10^6$ Riemann zeros (Odlyzko dataset) over
$\varepsilon \in [10^{-6}, 10^{-1}]$ must yield a stable linear region in
$\log N(\varepsilon)$ vs.\ $\log(1/\varepsilon)$ across at least two decades.
$D_H(n)$ must converge (not oscillate or diverge) across $n = 10^4, 10^5, 10^6$.
Failure to converge falsifies claim~(i).
\end{prediction}

\subsubsection*{What Would Falsify It}

If $D_H(n)$ oscillates persistently with no stable attractor, or if no scaling regime
appears over two decades of $\varepsilon$, Conjecture~B is falsified. Partial falsification:
if $D_H$ converges but lies outside $(1,2)$, claim~(ii) is falsified while claim~(i)
survives.

\bigskip

% ── Conjecture C ──────────────────────────────────────────────────────────────
\subsection{Conjecture C: Geometric Foundation}

\subsubsection*{Statement}

\begin{conjecture}[Geometric Foundation]
Every numerical mathematical relationship $R$ on a set $A \subseteq \mathbb{Z}$ has a
geometric realization: there exists a weighted graph $G = (V, E, w)$ with weight function
$\varphi: E \to A$ such that:
\begin{conjbox}
\[
  R \text{ holds in } A \;\iff\; R_G \text{ holds in } G,
\]
\end{conjbox}
where $R_G$ is the corresponding structural property of $G$. Furthermore, the geometric
realization is \emph{epistemically productive}: it exposes structural regularities via
spectral and topological methods inaccessible to coordinate-based numerical analysis of
$A$ directly.
\end{conjecture}

\subsubsection*{Derivation Sketch and Missing Step}

For Goldbach's conjecture: construct $G_n = (P, E_n)$ where $P$ is the set of primes
$\leq n$ and $\{p,q\} \in E_n$ iff $p + q = n$. Goldbach's conjecture is equivalent to:
for all even $n > 2$, $G_n$ has at least one edge. The spectral gap
$\lambda_2(G_n)$ (second Laplacian eigenvalue) measures graph connectivity~\cite{spielman2007}; if
$\lambda_2 > 0$ for all $n$, the graph is connected, implying the conjecture.

\textbf{Missing step:} The biconditional $R \iff R_G$ is established for specific cases.
The general claim---that \emph{every} numerical relationship has such a realization---requires
either a constructive proof for arbitrary $R$, or a categorical formalism showing the
embedding functor is fully faithful. Neither exists yet.

\subsubsection*{Falsifiable Predictions}

\begin{prediction}[PC]
The spectral gap $\lambda_2(G_n)$ for Goldbach graphs $G_n$,
$n \in \{4, 6, \ldots, 10^6\}$, should grow monotonically. Stagnation or decrease over
any run of 1000 consecutive even integers falsifies the epistemic-productivity claim.
Separately: geometric analysis (spectral gap trend) should predict Goldbach density faster
than direct numerical checking. Failure to do so falsifies epistemic productivity
independently of the biconditional.
\end{prediction}

\subsubsection*{What Would Falsify It}

Two independent paths: (1) $\lambda_2(G_n)$ fails monotonicity---falsifies epistemic
productivity. (2) A number-theoretic relationship is found for which no satisfying graph
realization exists---falsifies the existence claim. These are separable; (1) can be
falsified while (2) survives.

\bigskip
\hrule
\bigskip

% ── 4. Unified Implications ───────────────────────────────────────────────────
\section{Unified Implications}

\textbf{Riemann Hypothesis.} The First Law establishes that exhaustive zero-verification is
physically bounded. Conjecture~B reframes the target: instead of checking zeros
incrementally, compute $D_H$ as a finite-sample invariant. If $D_H$ is stable and tied to a
known geometric invariant of the zeta function, a finitary proof strategy becomes
conceivable---one that does not require the infinite-precision verification the First Law
rules out.

\textbf{P vs.\ NP.} Conjecture~A makes the frame-relativity of this question explicit. Any
claimed proof of $\mathbf{P} \neq \mathbf{NP}$ in the classical Turing model leaves open
whether $\mathbf{BQP} \supseteq \mathbf{NP}$ in a quantum substrate. The Second Law
provides the formal vocabulary for this distinction.

\textbf{Foundations.} Laws 1 and 2 together imply that realized mathematical truth is always
finite and approximate. The framework is compatible with Platonism about abstract objects. It
is incompatible with Platonism about \emph{realized} truth: no physical system can verify
infinite-precision statements, meaning the realized mathematical universe is strictly smaller
than the abstract one.

\bigskip
\hrule
\bigskip

% ── 5. Experimental Program ───────────────────────────────────────────────────
\section{Experimental Program}

Predictions P1, P2a, P2b test the Laws. PA, PB, PC test the Conjectures. All five are
independently executable with current or near-current technology.

\begin{enumerate}[label=\textbf{(\arabic*)}, leftmargin=*, itemsep=6pt]

  \item \textbf{P1 (First Law).} Compute $10^3, 10^4, 10^5$ Riemann zeros via
    \texttt{mpmath} on GPU. Monitor energy via hardware power interface. Fit
    $E(n) = \alpha \cdot n \log n + \beta$; verify $\alpha \geq k_B T \ln 2 \cdot c_0$.

  \item \textbf{P2a (Second Law).} Lorenz RK4, $p \in \{32, 64, 128\}$ bits, fixed $\tau$,
    pre-chaos regime. Fit $|\epsilon| = A \cdot 2^{-p} + B \cdot h^4$. Both terms must be
    measurably significant.

  \item \textbf{P2b (Second Law).} Identify empirical $h^*$ at each $p$; compare to analytic
    minimizer. Match required within experimental precision.

  \item \textbf{PA (Conjecture A).} Grover 3-SAT, 20 qubits, IBM Qiskit. Runtime vs.\
    $\sqrt{N}$ and $N$. Report $\gamma$ (decoherence rate) as covariate.
    Fit $t_q(N) = c \cdot \sqrt{N} \cdot e^{\gamma t}$.

  \item \textbf{PB (Conjecture B).} Box-counting of first $10^6$ Riemann zeros (Odlyzko).
    Plot $\log N(\varepsilon)$ vs.\ $\log(1/\varepsilon)$. Extract $D_H$ from linear region.
    Track $D_H(n)$ across $n = 10^4, 10^5, 10^6$ for convergence.

  \item \textbf{PC (Conjecture C).} NetworkX Goldbach graphs, $n \in [4, 10^6]$. Compute
    $\lambda_2$ via \texttt{scipy.sparse.linalg.eigsh}. Plot $\lambda_2(n)$ vs.\ $n$.
    Flag any monotonicity violations.

\end{enumerate}

\bigskip
\hrule
\bigskip

% ── 6. Future Directions ──────────────────────────────────────────────────────
\section{Future Directions}

\begin{itemize}[itemsep=6pt]
  \item \textbf{First Law.} Compute substrate-specific $n_{\max}$ for existing
    supercomputers, quantum annealers, and neuromorphic hardware. Publish as a physical
    precision budget table.

  \item \textbf{Second Law.} Formalize inter-frame error accumulation for distributed
    computation. Does network latency between nodes constitute a frame boundary? Derive the
    distributed analogue of the two-term error bound.

  \item \textbf{Conjecture A $\to$ Law.} Characterize which physical substrates instantiate
    $\mathbf{BQP}$ formally, from thermodynamic constraints rather than algorithm design.
    If successful, Conjecture~A becomes the Third Law.

  \item \textbf{Conjecture B $\to$ Law.} Prove $D_H(n)$ convergence from GUE
    scale-invariance rigorously. Test whether $D_H$ is conserved across the $L$-function
    family. If conserved, $D_H$ is a candidate for a new arithmetic invariant.

  \item \textbf{Conjecture C $\to$ Law.} Develop the categorical formalism for the embedding
    $\varphi: E \to A$. Prove or disprove that the functor is fully faithful for all
    number-theoretic $R$.
\end{itemize}

\bigskip
\hrule
\bigskip

% ── 7. Conclusion ─────────────────────────────────────────────────────────────
\section{Conclusion}

Anson's Framework for Physical Mathematics relocates mathematical truth: not from the
Platonic realm of timeless abstraction to something lesser, but to something more
honest---the physical realm of constrained, frame-dependent, energetically bounded
realization. Two laws are established from first principles and carry the full epistemic
weight of that word. Three conjectures are held to the same formal standard and are waiting
on proof steps and experimental confirmation to join them.

The distinction between laws and conjectures is not a hedge. It is the posture of a
framework that intends to survive contact with evidence. Each component is falsifiable. The
conjectures are not speculations---they are the next tier of the same structure, already
dressed for the occasion.

\bigskip
\hrule
\bigskip

% ── Author and Process ────────────────────────────────────────────────────────
\section*{Author and Process}

\textbf{Author.} Amber Anson, independent researcher, AmberContinuum Research.

\textbf{Process.} The ideas, conceptual architecture, and epistemological positions are
the original work of the author. The mathematical formalization was developed through
human-AI coupling with Claude (Anthropic) and Grok~3 (xAI), functioning as coherence
engines and formal reasoning substrates. These systems did not originate the claims; they
made the author's claims mathematically legible.

The author regards this paper as a documented instance of the coupling dynamics described
by the Perturbation Dynamics Operator Theory (PDOT) framework
(pre-registered OSF: osf.io/7mejx). The Lean~4 formalization is the proof of record.
The process that produced it is the proof of concept.

\bigskip
\hrule
\bigskip

% ── References ────────────────────────────────────────────────────────────────
\begin{thebibliography}{99}

\bibitem{lean4}
de Moura, L., \& Ullrich, S. (2021).
The Lean~4 theorem prover and programming language.
In \textit{Automated Deduction -- CADE 28}, Lecture Notes in Computer Science, Vol. 12699 (pp. 625--635). Springer.
\url{https://doi.org/10.1007/978-3-030-79876-5_37}

\bibitem{mathlib2020}
The mathlib Community. (2020).
The Lean mathematical library.
In \textit{Proceedings of the 9th ACM SIGPLAN International Conference on Certified Programs and Proofs (CPP 2020)} (pp. 367--381). ACM.
\url{https://doi.org/10.1145/3372885.3373824}

\bibitem{bekenstein1981}
Bekenstein, J.~D. (1981).
Universal upper bound on the entropy-to-energy ratio for bounded systems.
\textit{Physical Review D}, 23(2), 287.

\bibitem{berry1999}
Berry, M.~V., \& Keating, J.~P. (1999).
The Riemann zeros and eigenvalue asymptotics.
\textit{SIAM Review}, 41(2), 236--266.

\bibitem{deutsch1985}
Deutsch, D. (1985).
Quantum theory, the Church--Turing principle and the universal quantum computer.
\textit{Proceedings of the Royal Society A}, 400, 97--117.

\bibitem{falconer2003}
Falconer, K. (2003).
\textit{Fractal Geometry: Mathematical Foundations and Applications}.
Wiley.

\bibitem{grover1996}
Grover, L.~K. (1996).
A fast quantum mechanical algorithm for database search.
\textit{Proceedings of the 28th ACM STOC}, 212--219.

\bibitem{higham2002}
Higham, N.~J. (2002).
\textit{Accuracy and Stability of Numerical Algorithms}.
SIAM.

\bibitem{johansson2023}
Johansson, F., et al. (2023).
mpmath: A Python library for arbitrary-precision floating-point arithmetic.
\url{http://mpmath.org}.

\bibitem{landauer1961}
Landauer, R. (1961).
Irreversibility and heat generation in the computing process.
\textit{IBM Journal of Research and Development}, 5(3), 183--191.

\bibitem{montgomery1973}
Montgomery, H.~L. (1973).
The pair correlation of zeros of the zeta function.
In H.~G. Diamond (Ed.), \textit{Analytic Number Theory}, Proc. Sympos. Pure Math., Vol.~XXIV (pp. 181--193). American Mathematical Society.

\bibitem{odlyzko1987}
Odlyzko, A.~M. (1987).
On the distribution of spacings between zeros of the zeta function.
\textit{Mathematics of Computation}, 48(177), 273--308.

\bibitem{shor1997}
Shor, P.~W. (1997).
Polynomial-time algorithms for prime factorization and discrete logarithms on a quantum computer.
\textit{SIAM Journal on Computing}, 26(5), 1484--1509.

\bibitem{spielman2007}
Spielman, D.~A. (2007).
Spectral graph theory and its applications.
\textit{Proceedings of the 48th Annual IEEE Symposium on Foundations of Computer Science (FOCS)}, 29--38.

\bibitem{stein1964}
Stein, M.~L., \& Ulam, S.~M. (1964).
An observation on the distribution of primes.
\textit{American Mathematical Monthly}, 71(5), 496--498.

\bibitem{virtanen2020}
Virtanen, P., et al. (2020).
SciPy 1.0: Fundamental algorithms for scientific computing in Python.
\textit{Nature Methods}, 17, 261--272.

\bibitem{wheeler1990}
Wheeler, J.~A. (1990).
Information, physics, quantum: The search for links.
In W.~H. Zurek (Ed.), \textit{Complexity, Entropy, and the Physics of Information} (pp. 309--336). Addison-Wesley, Redwood City, CA.

\end{thebibliography}

\bigskip
\hrule
\bigskip

% ── Appendix: Lean 4 Formal Verification ─────────────────────────────────────
\appendix
\section{Lean 4 Formal Verification}

The following file contains the complete Lean~4~\cite{lean4} (Mathlib~\cite{mathlib2020}) formalization of
Anson's First and Second Laws, verified mechanically by the Lean type checker.
The file can be pasted directly into the Lean~4 web editor at
\url{https://live.lean-lang.org} with Mathlib selected to reproduce the
verification. The three research claims are recorded using \texttt{def ... : Prop}
declarations. Their current formal statements have counterexamples, so Lean
neither assumes nor proves them; each needs revised hypotheses or encoding.

\bigskip

\begingroup
\small
\begin{verbatim}
-- Physical Mathematics: A Foundational Framework
-- Formal Verification of Anson's First and Second Laws
-- Amber Anson, AmberContinuum Research

import Mathlib

set_option linter.unusedVariables false

namespace PhysicalMathematics

open Real

variable {k_B T R E hbar c : R}

-- ===========================================================
-- PRIMITIVE DEFINITIONS
-- ===========================================================

-- Landauer cost: minimum energy per irreversible bit operation
-- E_bit = k_B * T * ln 2
noncomputable def landauerCost (k_B T : R) : R :=
  k_B * T * log 2

-- Bekenstein bound: maximum entropy of a bounded physical system
-- S_max = (2pi * R * E) / (hbar * c)
noncomputable def bekensteinBound (R E hbar c : R) : R :=
  (2 * pi * R * E) / (hbar * c)

-- Maximum verifiable bits
-- n_max = floor( (2pi R E) / (hbar c k_B T ln 2) )
noncomputable def nMax (R E hbar c k_B T : R) : N :=
  floor(bekensteinBound R E hbar c / (k_B * T * log 2))_+

-- ===========================================================
-- LEMMAS
-- ===========================================================

lemma log_two_pos : 0 < log 2 := by apply log_pos; norm_num

lemma landauerCost_pos (hk_B : 0 < k_B) (hT : 0 < T) :
    0 < landauerCost k_B T :=
  mul_pos (mul_pos hk_B hT) log_two_pos

lemma bekensteinBound_pos
    (hR : 0 < R) (hE : 0 < E) (hhbar : 0 < hbar) (hc : 0 < c) :
    0 < bekensteinBound R E hbar c := by
  unfold bekensteinBound
  apply div_pos _ (mul_pos hhbar hc)
  have : (0 : R) < 2 * pi * R * E := by positivity
  linarith

-- ===========================================================
-- ANSON'S FIRST LAW: THERMODYNAMIC CONSTRAINT
-- ===========================================================

theorem law1_ratio_pos
    (hR : 0 < R) (hE : 0 < E) (hhbar : 0 < hbar) (hc : 0 < c)
    (hk_B : 0 < k_B) (hT : 0 < T) :
    0 < bekensteinBound R E hbar c / (k_B * T * log 2) :=
  div_pos (bekensteinBound_pos hR hE hhbar hc)
          (mul_pos (mul_pos hk_B hT) log_two_pos)

theorem law1_energy_ceiling
    (hk_B : 0 < k_B) (hT : 0 < T) (hE : 0 < E)
    (b : N) (hb : E < b * (k_B * T * log 2)) :
    not (b * (k_B * T * log 2) <= E) :=
  not_le.mpr hb

theorem law1_bekenstein_ceiling
    (hR : 0 < R) (hE : 0 < E) (hhbar : 0 < hbar) (hc : 0 < c)
    (b : N) (hb : bekensteinBound R E hbar c < b * log 2) :
    not (b * log 2 <= bekensteinBound R E hbar c) :=
  not_le.mpr hb

-- Core theorem: any computation requiring more than n_max bits
-- exceeds the physical entropy budget of the system.
theorem law1_core
    (hR : 0 < R) (hE : 0 < E) (hhbar : 0 < hbar) (hc : 0 < c)
    (hk_B : 0 < k_B) (hT : 0 < T)
    (b : N) (hb : nMax R E hbar c k_B T < b) :
    bekensteinBound R E hbar c < b * (k_B * T * log 2) := by
  unfold nMax at hb
  have hdenom := mul_pos (mul_pos hk_B hT) log_two_pos
  have hBpos := bekensteinBound_pos hR hE hhbar hc
  have hpos  := div_pos hBpos hdenom
  have hfloor := Nat.floor_le (le_of_lt hpos)
  have hb' : (floor(...)+  : R) < (b : R) := by exact_mod_cast hb
  have hle : bekensteinBound R E hbar c / (k_B * T * log 2) < (b : R) := by
    have hstep := Nat.lt_floor_add_one (bekensteinBound R E hbar c / ...)
    have hcast : (floor(...)+  : R) + 1 <= (b : R) := by exact_mod_cast hb
    linarith
  rw [div_lt_iff_0 hdenom] at hle
  linarith

-- ===========================================================
-- ANSON'S SECOND LAW: COMPUTATIONAL FRAME DEPENDENCE
-- ===========================================================

variable {C_method C_round h p : R} {q : N}

lemma truncation_pos (hCm : 0 < C_method) (hh : 0 < h) (q : N) :
    0 < C_method * h ^ q :=
  mul_pos hCm (pow_pos hh q)

lemma rounding_pos (hCr : 0 < C_round) (hp : 0 < p) :
    0 < C_round * p^(-1) :=
  mul_pos hCr (inv_pos.mpr hp)

-- (i) Both error terms are strictly positive in any finite frame.
theorem law2_error_bound_pos
    (hCm : 0 < C_method) (hCr : 0 < C_round)
    (hh : 0 < h) (hp : 0 < p) (q : N) :
    0 < C_method * h ^ q + C_round * p^(-1) :=
  add_pos (truncation_pos hCm hh q) (rounding_pos hCr hp)

-- (ii-a) Truncation error is strictly increasing in h.
theorem law2_truncation_monotone
    (hCm : 0 < C_method) (hh_1 : 0 < h_1)
    (hh_ord : h_1 < h_2) (q : N) (hq : 0 < q) :
    C_method * h_1 ^ q < C_method * h_2 ^ q :=
  mul_lt_mul_of_pos_left
    (pow_lt_pow_left_0 hh_ord (le_of_lt hh_1)
      (Nat.pos_iff_ne_zero.mp hq)) hCm

-- (ii-b) Rounding error is strictly decreasing in p.
theorem law2_rounding_monotone
    (hCr : 0 < C_round) (hp_1 : 0 < p_1) (hp_2 : 0 < p_2)
    (hp_ord : p_1 < p_2) :
    C_round * p_2^(-1) < C_round * p_1^(-1) := by
  apply mul_lt_mul_of_pos_left _ hCr
  have hp1inv : p_1^(-1) = 1 / p_1 := (one_div p_1).symm
  have hp2inv : p_2^(-1) = 1 / p_2 := (one_div p_2).symm
  rw [hp1inv, hp2inv]
  have := div_lt_div_of_pos_left (by norm_num : (0:R) < 1) hp_1 hp_ord
  linarith [this]

-- (ii) Tradeoff: both monotonicity results as a conjunction.
theorem law2_tradeoff
    (hCm : 0 < C_method) (hCr : 0 < C_round)
    (hh_1 : 0 < h_1) (hp_1 : 0 < p_1) (hp_2 : 0 < p_2)
    (hh_ord : h_1 < h_2) (hp_ord : p_1 < p_2)
    (q : N) (hq : 0 < q) :
    C_method * h_1 ^ q < C_method * h_2 ^ q /\
    C_round * p_2^(-1) < C_round * p_1^(-1) :=
  And.intro
    (law2_truncation_monotone hCm hh_1 hh_ord q hq)
    (law2_rounding_monotone hCr hp_1 hp_2 hp_ord)

-- (iii) Optimal step is well-defined and strictly positive.
noncomputable def optimalStep (C_method C_round p : R) (q : N) : R :=
  (C_round / (q * C_method * p)) ^ ((1 : R) / (q + 1))

theorem optimalStep_pos
    (hCm : 0 < C_method) (hCr : 0 < C_round)
    (hp : 0 < p) (q : N) (hq : 0 < q) :
    0 < optimalStep C_method C_round p q := by
  unfold optimalStep
  apply rpow_pos_of_pos
  apply div_pos hCr
  apply mul_pos; apply mul_pos
  · exact_mod_cast hq
  · exact hCm
  · exact hp

-- ===========================================================
-- COUPLING OF LAWS 1 AND 2
-- ===========================================================

-- The precision p in Law 2 is bounded by n_max from Law 1.
theorem laws_coupled
    (hR : 0 < R) (hE : 0 < E) (hhbar : 0 < hbar) (hc : 0 < c)
    (hk_B : 0 < k_B) (hT : 0 < T)
    (p : N) (hp : nMax R E hbar c k_B T < p) :
    bekensteinBound R E hbar c < p * (k_B * T * log 2) :=
  law1_core hR hE hhbar hc hk_B hT p hp

-- ===========================================================
-- RESEARCH PROPOSITIONS (unasserted; current forms need revision)
-- ===========================================================

-- Conjecture A: Complexity class membership is substrate-dependent.
def conjectureA_quantum_complexity_reduction
    (C_classical C_quantum : N -> R)
    : Prop :=
    (forall n, 0 < C_classical n) ->
    (forall n, 0 < C_quantum n) ->
    forall n, C_quantum n < C_classical n

-- Conjecture B: Hausdorff dimension of a number-theoretic
-- sequence converges to a stable limit in (1, 2).
def conjectureB_fractal_dimension_convergence
    (D : N -> R) : Prop :=
    (forall n, 1 < D n /\ D n < 2) ->
    exists (D_H : R), 1 < D_H /\ D_H < 2 /\
      Filter.Tendsto D Filter.atTop (nhds D_H)

-- Conjecture C: Every numerical relationship has a geometric
-- realization via a graph embedding.
def conjectureC_geometric_foundation
    (R : Z -> Z -> Prop) (R_G : (Z -> Z -> Prop) -> Prop) : Prop :=
    exists (phi : (Z -> Z -> Prop) -> (Z -> Z -> Prop)),
      forall (a b : Z), R a b <-> R_G (phi R)

end PhysicalMathematics
\end{verbatim}
\endgroup

\bigskip
\noindent\textbf{Note on verbatim rendering.} The appendix above uses ASCII
approximations for certain Unicode symbols used in the actual Lean~4 source
(e.g., \texttt{hbar} for $\hbar$, \texttt{floor(...)\_+} for $\lfloor\cdot\rfloor_+$,
\texttt{p\^{}(-1)} for $p^{-1}$, \texttt{/\textbackslash} for $\wedge$).
The canonical source file \texttt{PhysicalMathematics.lean} contains the full
Unicode Lean~4 syntax and is the authoritative verification artifact.

\end{document}
